% ============================================================
% VI. CONCLUSIONES
% ============================================================

\section{Conclusiones}

Se diseñaron, simularon e implementaron filtros digitales FIR e IIR cumpliendo especificaciones de frecuencia de corte, ganancia de banda rechazada y ancho de banda de transición dadas, confirmando teóricamente que el IIR logra las mismas especificaciones de magnitud con un orden —y por tanto un costo computacional y de memoria— sensiblemente menor que el FIR equivalente, a cambio de una fase no lineal que se manifiesta como un desfase dependiente de la frecuencia, frente al desfase proporcional a la frecuencia (fase lineal) propio del FIR; ambos efectos fueron caracterizados cuantitativamente en las Tablas~\ref{tab:ganancia} y \ref{tab:desfase}.

Para el filtrado de las señales reales de sensores del robot (IMU y encoders), el análisis frecuencial permitió separar claramente la banda de la señal de interés (dinámica lenta del movimiento del robot, $0$--$4$ Hz) de la banda del ruido (banda ancha, por encima de $10$--$15$ Hz), lo que favoreció un filtro IIR convencional de orden bajo como mejor equilibrio entre atenuación de ruido lograda y costo computacional frente a alternativas de media móvil y autorregresivas, ambas también efectivas pero sin ventaja clara sobre el IIR en este caso particular, donde el ruido es razonablemente estacionario en frecuencia.

La integración final —un robot Minisumo diferencial teleoperado mediante ROS2, con control PID de velocidad por motor y odometría calculada a partir de las velocidades de encoder filtradas— mostró, incluso en el análisis puramente teórico presentado en este documento, que el error de posición estimada por odometría frente a la posición real crece con la distancia recorrida de cada trayectoria, una limitación inherente al método de odometría de lazo abierto y no al filtrado aplicado, el cual sí mejora la calidad de la velocidad integrada pero no elimina el deslizamiento de ruedas ni el error de calibración geométrica del robot.

\textbf{Este informe es una plantilla ilustrativa}: ningún resultado proviene de mediciones físicas; los valores de las Tablas~\ref{tab:ganancia}--\ref{tab:posicion} son ilustrativos, incluidos para mostrar la estructura esperada del informe final. El trabajo pendiente directo es: ensamblar el robot Minisumo con su IMU y encoders, implementar los filtros diseñados en el microcontrolador, capturar datos reales con el generador de ondas y el osciloscopio para las señales periódicas, y con el robot en movimiento dentro del dojo para las señales de sensores y la odometría, reemplazando cada valor ilustrativo por su medición real y verificando las tendencias cualitativas predichas en la Discusión.

Como trabajo futuro se sugiere evaluar un filtro adaptativo propiamente dicho (p. ej. LMS) bajo condiciones de ruido no estacionario (distintas superficies del dojo, distintas velocidades de motor), y complementar la odometría con una fuente de corrección externa (p. ej. visión desde una cámara fija sobre el dojo) para acotar el crecimiento del error de posición en trayectorias largas.
